\chapter{Hybrid Dynamics and Other Topics}
\label{chap:9}
This chapter covers the following topics: hybrid dynamics, floating bases, gears and dynamic equivalence. Hybrid dynamics is a generalization of forward and inverse dynamics in which the forces are known at some joints, the accelerations at the rest, and the task is to calculate the unknown forces and accelerations. In effect, one is performing forward dynamics at some joints and inverse dynamics at the rest. A floating-base system is one in which the base is a moving body. Such systems can be modelled as fixed-base systems by installing a 6DoF joint between the fixed base and the body representing the floating base. However, floating-base systems are an important class of rigid-body systems, and merit special treatment. Sections 9.1 to 9.5 present several algorithms for hybrid dynamics and floating bases. Gears impose constraints between joint variables. They therefore resemble kinematic loops, and Section 9.6 shows how gear constraints can be incorporated into a rigid-body system using a technique that was developed in Chapter 8 for kinematic loops. Two different rigid-body systems are dynamically equivalent if they have the same equation of motion. Section 9.7 explores this concept, and presents a method for modifying the inertia parameters of a rigid-body system without altering its equation of motion. 9.1 Hybrid Dynamics For each joint i, there is a force variable, τi, and an acceleration variable, ¨qi. In the forward-dynamics problem, we are given τi for every joint, and the task is to calculate the accelerations. In the inverse-dynamics problem, we are given ¨qi for every joint, and the task is to calculate the forces. In the hybriddynamics problem, we are given either τi or ¨qi at each joint, and the task is to calculate the unknown accelerations and forces. In effect, hybrid dynamics is like performing forward dynamics at some of the joints and inverse dynamics

at the rest. The hybrid-dynamics problem therefore includes both the forwarddynamics problem and the inverse-dynamics problem as special cases. Hybrid dynamics is used to introduce prescribed motions into a rigid-body system. For example, it will let you study the dynamics of a mechanical linkage or transmission under prescribed motion conditions at the input, or the motion of a compound pendulum in response to a given kinematic excitation, or the tumbling behaviour of an animated character in mid air. It can also be used to help design trajectories for underactuated systems; for example, a sequence of limb movements for an animated character to leap, perform a somersault, and land on its feet. Another use is to simplify a simulation by removing unimportant and/or high-frequency dynamics. For example, the motions of a robot’s joints are typically controlled by high-gain feedback control systems. These control systems continually compare the actual motion of a joint with the commanded motion, and adjust the force output from the actuator so as to bring the difference quickly to zero. In a simulation of a complete robot, the simulator must calculate the dynamics of the actuators and control systems as well as the mechanical parts. However, if a particular control system can be regarded as perfect, meaning that the actual motions of the joints it controls are practically equal to the commanded motions, then it is possible to dispense with this control system and its actuators, and simply prescribe the motions of the affected joints. If the removed components were responsible for the highest-frequency dynamics in the original system, then the simulator can now use a longer integration time step. An example of this technique can be found in Hu et al. (2005). Equations Let us define fd to be the set of ‘forward-dynamics joints’; that is, the set of joints for which the forces are known and the accelerations unknown. Any joint that is not a member of fd is therefore an ‘inverse-dynamics joint’, for which the acceleration is known and the force unknown. Let us also define Q to be the permutation matrix that reorders the elements of ¨q and τ so that the forward-dynamics variables come first. Thus, Q ¨q = ¨q1 ¨q2  and Q τ = τ1 τ2  , where ¨q1 and τ1 contain the acceleration and force variables pertaining to the forward-dynamics joints, and ¨q2 and τ2 contain the acceleration and force variables pertaining to the inverse-dynamics joints. The equation of motion for a kinematic tree can now be written  H11 H12 H21 H22  ¨q1 ¨q2  =  τ1 τ2  −  C1 C2  , (9.1) where H11 H12 H21 H22  = Q H QT and C1 C2  = Q C .

Equation 9.1 is just a permutation of the standard equation of motion, H ¨q = τ −C (e.g. see Eq. 6.1). The two unknowns in this equation are ¨q1 and τ2. If we bring them together on the left-hand side, then we get  H11 0 H21 −1  ¨q1 τ2  =  τ1 0  −  C′ 1 C′ 2  , (9.2) where C′ 1 C′ 2  = C1 + H12¨q2 C2 + H22¨q2  = Q C′ . This is the basic equation for hybrid dynamics. Physically, C′ is the force required to impart zero acceleration to each forward-dynamics joint and the given acceleration to each inverse-dynamics joint. Thus, C′ can be calculated by an inverse-dynamics function as follows: C′ = ID(q, ˙q, QT  0 ¨q2  ) . (9.3) Example 9.1 Equation 9.2 hides ¨q2 within C′. If we want to keep ¨q2 visible, then the equation can be written H11 0 H21 −1  ¨q1 τ2  = 1 −H12 0 −H22  τ1 ¨q2  − C1 C2  . (9.4) Multiplying this equation by the inverse of the left-hand coefficient matrix produces Eq. 10 in Jain and Rodriguez (1993). Algorithms Equation 9.2 can be solved in four steps as follows. 1. Calculate C′ using Eq. 9.3. 2. Calculate H11 as described below. 3. Solve H11¨q1 = τ1 −C′ 1 for ¨q1. 4. Calculate τ2 via τ = C′ + IDδ(QT ¨q1 0  ). The symbol IDδ in step 4 denotes the differential inverse-dynamics function described in Table 6.1, which multiplies its argument by H. Thus, the equation in step 4 evaluates to τ = C′ + QT H11¨q1 H21¨q1  . The simplest way to accomplish step 2 is to calculate H via the composite rigid body algorithm (§6.2) and extract the submatrix H11. However, if greater

H = 0 for each i in ν(fd) do Ic i = Ii end for i = NB to 1 do if λ(i) ∈ν(fd) then Ic λ(i) = Ic λ(i) + λ(i)X∗ i Ic i iXλ(i) end if i ∈fd then F = Ic i Si Hii = ST i F ... ... j = i while λ(j) ∈ν(fd) do F = λ(j)X∗ j F j = λ(j) if j ∈fd then Hij = F TSj Hji = HT ij end end end end Table 9.1: Modified composite rigid body algorithm for calculating H11 efficiency is required, then it is possible to modify the algorithm so that it only performs the minimum amount of calculation necessary to obtain H11. This modified algorithm is shown in Table 9.1. One new quantity appears in this algorithm: ν(fd), which is the set of bodies that are supported by at least one forward-dynamics joint. It is defined by the formula ν(fd) = [ i∈fd ν(i) , (9.5) and it can be calculated from fd and λ as follows:1 ν(fd) = fd for i = 1 to NB do if λ(i) ∈ν(fd) then ν(fd) = ν(fd) ∪\{i\} end end The cost of calculating ν(fd) is small enough that it would be practical for a hybrid dynamics routine to calculate it each time the routine was called. However, as ν(fd) depends only on fd and λ, its value would not normally change during the course of a single simulation run; so it would also be possible to calculate it just once at the beginning of the run, and store it for later use. On comparing the modified algorithm with the original, in Table 6.2, it can be seen that they both have the same structure, but the modified algorithm performs only a subset of the calculations. For example, I c i is calculated only 1In this code fragment, ‘ν(fd)’ should be regarded as the name of a variable. In a practical implementation, ν(fd) could be represented by an array of NB + 1 Boolean values, the array being indexed from 0 to NB, such that each element is either true or false according to whether its index is or is not an element of ν(fd).

inverse dynamics joints 1 2 3 4 λ = [0,1,1,2,2,3,3] λ’ = [0,0,0,1,1] 5 6 7 0 0 1 2 3 4 5 H11 H original graph modified graph H11 Calculation of λ′ map(0) = 0 j = 1 for i = 1 to n do if i ∈fd then map(i) = j λ′(j) = map(λ(i)) j = j + 1 else map(i) = map(λ(i)) end end Figure 9.1: Branch-induced sparsity in H11 for those values of i satisfying i ∈ν(fd), and Hij is calculated only for those values of i and j satisfying i, j ∈fd. Moving on to step 3, H11 is a symmetric, positive-definite matrix, so the equation in step 3 could be solved using any standard linear-equation solver for such matrices. However, as shown in Figure 9.1, H11 inherits a portion of whatever branch-induced sparsity is present in H. We therefore have the opportunity to exploit this sparsity by using one of the algorithms described in Section 6.5. The prerequisite for exploiting the sparsity in H11 is to obtain a modified parent array, λ′, that describes the sparsity pattern in this matrix. Figure 9.1 shows how this can be done. Starting with the connectivity graph of the original system, a modified graph is constructed by merging every pair of nodes that are connected by an inverse-dynamics joint. Once this is done, the remaining nodes must be renumbered, as there are now fewer of them. In the example shown in the figure, joints 1 and 2 are inverse-dynamics joints, so nodes 1 and 2 are merged with node 0, and nodes 3 to 7 are renumbered 1 to 5. Figure 9.1 also shows an algorithm for calculating λ′ from λ and fd. As a by-product, this algorithm also calculates the node-number mapping from the original to the modified graph: map(i) is the new number for original node i.

To follow this algorithm, it helps to know that i is counting node numbers in the original graph, while j is counting node numbers in the modified graph. Bear in mind that if any of the joints in the modified graph have more than one degree of freedom then λ′ will have to be expanded, as explained in Section 6.5, before it can be used. 9.2 Articulated-Body Hybrid Dynamics The articulated-body algorithm is easily adapted to perform hybrid dynamics: one simply modifies the algorithm so that it uses the normal articulated-body equations (as described in §7.3) at each forward-dynamics joint, and the equations developed below at each inverse-dynamics joint. The two sets of equations are shown side-by-side in Table 9.2, and the pseudocode is shown in Table 9.3. The latter is adapted from the code in Table 7.1. Let fi be the spatial force transmitted across joint i. fi is related to the acceleration of body i by the equation fi = IA i ai + pA i ; (9.6) and it is related to the acceleration of body λ(i) by the equation fi = Ia i aλ(i) + pa i . (9.7) (See Eq. 7.25.) Furthermore, aλ(i) is related to ai by ai = aλ(i) + ˙Si ˙qi + Si ¨qi . (9.8) In the articulated-body algorithm, the main step is the calculation of I a i and pa i from IA i and pA i . If joint i is a forward-dynamics joint, then the immediate problem is that ¨qi is unknown, so the first step is to find an expression for ¨qi. However, no such difficulty exists for an inverse-dynamics joint. Substituting Eq. 9.8 into Eq. 9.6 gives fi = IA i (aλ(i) + ˙Si ˙qi + Si ¨qi) + pA i ; and comparing this equation with Eq. 9.7 reveals that Ia i = IA i (9.9) and pa i = pA i + IA i ( ˙Si ˙qi + Si ¨qi) . (9.10) These two equations take the place of Eqs. 7.39 and 7.40 in Pass 2. Apart from this, only two other changes can be seen in the inverse-dynamics equations in Table 9.2: the term Si¨qi has been added to ci (which means that a′ i no longer appears in Pass 3), and a new equation has been added to Pass 3 to calculate τi via τi = ST i fi and Eq. 9.6. (We could have used Eq. 9.7 instead.)

Forward Dynamics (i ∈fd): Pass 1 vJi = Si ˙qi vi = iXλ(i) vλ(i) + vJi (v0 = 0) cJi = ˚ Si ˙qi ci = cJi + vi × vJi pi = vi ×∗Ii vi −iX∗ 0 f x i Inverse Dynamics (i /∈fd): Pass 1 vJi = Si ˙qi vi = iXλ(i) vλ(i) + vJi (v0 = 0) cJi = ˚ Si ˙qi ci = cJi + vi × vJi + Si ¨qi pi = vi ×∗Ii vi −iX∗ 0 f x i Pass 2 IA i = Ii + X j∈µ(i) iX∗ j Ia j jXi pA i = pi + X j∈µ(i) iX∗ j pa j Ui = IA i Si Di = ST i Ui ui = τi −ST i pA i Ia i = IA i −Ui D−1 i U T i pa i = pA i + Ia i ci + Ui D−1 i ui Pass 2 IA i = Ii + X j∈µ(i) iX∗ j Ia j jXi pA i = pi + X j∈µ(i) iX∗ j pa j Ia i = IA i pa i = pA i + Ia i ci Pass 3 a′ i = iXλ(i) aλ(i) + ci (a0 = −ag) ¨qi = D−1 i (ui −U T i a′ i) ai = a′ i + Si ¨qi Pass 3 ai = iXλ(i) aλ(i) + ci (a0 = −ag) τi = ST i (IA i ai + pA i ) Table 9.2: Articulated-body hybrid dynamics equations

v0 = 0 for i = 1 to NB do [XJ, Si, vJ, cJ] = jcalc(jtype(i), qi, ˙qi) iXλ(i) = XJ XT(i) if λ(i)̸ = 0 then iX0 = iXλ(i) λ(i)X0 end vi = iXλ(i) vλ(i) + vJ if i ∈fd then ci = cJ + vi × vJ else ci = cJ + vi × vJ + Si ¨qi end IA i = Ii pA i = vi ×∗Ii vi −iX∗ 0 f x i end for i = NB to 1 do if i ∈fd then Ui = IA i Si Di = ST i Ui ui = τi −ST i pA i if λ(i)̸ = 0 then Ia = IA i −Ui D−1 i U T i IA λ(i) = IA λ(i) + λ(i)X∗ i Ia iXλ(i) pa = pA i + Iaci + Ui D−1 i ui pA λ(i) = pA λ(i) + λ(i)X∗ i pa end else if λ(i)̸ = 0 then Ia = IA i IA λ(i) = IA λ(i) + λ(i)X∗ i Ia iXλ(i) pa = pA i + Iaci pA λ(i) = pA λ(i) + λ(i)X∗ i pa end end end ... ... a0 = −ag for i = 1 to NB do if i ∈fd then a′ = iXλ(i) aλ(i) + ci ¨qi = D−1 i (ui −U T i a′) ai = a′ + Si ¨qi else ai = iXλ(i) aλ(i) + ci τi = ST i (IA i ai + pA i ) end end Table 9.3: Articulated-body hybrid dynamics algorithm

\section{loating Bases A floating-base rigid-body system is one in which the base body (body 0) is free to move, rather than being fixed in space. Any floating-base system can be converted into an equivalent fixed-base system by making the following changes: 1. Add 1 to all the body and joint numbers, so that the floating base is now body 1, and the lowest-numbered joint is joint 2. 2. Add a fixed base to the system, and a new 6-DoF joint connected between the fixed and floating bases; the former being the new body 0, and the latter the new joint 1. 3. As a consequence of step 2, add 1 to NB and NJ, and add 6 to n. Having made these changes, the floating-base system’s dynamics can be analysed and calculated by any of the techniques and algorithms designed for fixedbase systems. Nevertheless, floating-base systems are an important special case, and they do merit special treatment. Let us develop the equations of motion for a floating-base kinematic tree. The quickest approach is to start with the equation of motion for the equivalent fixed-base system. This equation can be written H11 H1∗ H∗1 H∗∗  ¨q1 ¨q∗  + C1 C∗  = τ1 τ∗  , (9.11) where the subscripts 1 and ∗refer to joint 1 and ‘all other joints’, respectively. This equation is already the correct equation of motion for a floating-base system—it just needs to be massaged into more appropriate notation, which is what we shall now do. The special property of joint 1 is that is has six degrees of freedom. Its motion-subspace matrix, S1, must therefore be a 6 × 6 full-rank matrix. In fact, S1 could be any 6 × 6 full-rank matrix, as two different choices of S1 amount only to two different choices of velocity variables, rather than to any fundamental change in the equation. We therefore choose the most convenient value, which is the identity matrix. In particular, we choose S1 to be the identity matrix in floating-base coordinates.2 Having specified S1 = 16×6, several other quantities now take on special values. Starting with ˙q1, ¨q1 and τ1, we have v1 = S1 ˙q1 = 16×6 ˙q1 = ˙q1 , a1 = ˙S1 ˙q1 + S1 ¨q1 = v1 × S1 ˙q1 + S1 ¨q1 = v1 × v1 + 16×6 ¨q1 = ¨q1 and τ1 = ST 1 f1 = 16×6 f1 = f1 . 2That is, 1S1 = 16×6; but in any other coordinate system we have iS1 = iX116×6, which is not the identity matrix. As S1 is fixed in body 1, we have ˚ S1 = 0 and ˙S1 = v1 × S1.}
\label{sec:9.3}
So ˙q1, ¨q1 and τ1 contain the Pl¨ucker coordinates of the spatial vectors v1, a1 and f1. Next, we observe that f1 and f x 1 have essentially the same meaning, since they are both external forces acting on the floating base; so the total external force is actually given by the sum f1 + f x 1 . In the equations below, f1 (and therefore also τ1) has been set to zero; but we always have the freedom to choose how to apportion the external force between f x 1 and f1. As C is the generalized bias force, it follows that C1 is the value that τ1 would have to take in order to produce zero joint acceleration at the floating base, assuming that all the other joint accelerations are zero. Given that ¨q1 = a1 and τ1 = f1, this means that C1 is also the spatial bias force for the composite rigid body comprising the whole floating-base system. We shall use the symbol pc 1 to denote this quantity. Finally, from the definition of Hij in Eq. 6.14, we have H1i = ST 1 Ic i Si = Ic i Si and H11 = ST 1 Ic 1 S1 = Ic 1 . Taking these changes into account, Eq. 9.11 can be written  Ic 1 H1∗ H∗1 H∗∗  a1 ¨q∗  + pc 1 C∗  =  0 τ∗  . Let us now revert to the original floating-base system, and adjust the notation accordingly. First, NB, NJ and n revert to their original values, and the bodies and joints revert to their original numbers. The symbols a1, pc 1, Ic 1, and so on, therefore become a0, pc 0, Ic 0, and so on. Furthermore, as the 6-DoF joint 1 no longer exists, the symbols H∗∗, C∗, ¨q∗and τ∗all simplify to H, C, ¨q and τ, respectively. We replace H1∗with a new symbol, F , which is defined as follows: F = F1 F2 · · · FNB  (9.12) where Fi = Ic i Si . F is a 6×n matrix whose columns are the spatial forces required at the floating base to support unit accelerations about each joint variable (cf. §6.3). With these changes, the equation of motion for a floating-base system now reads  Ic 0 F F T H  a0 ¨q  + pc 0 C  = 0 τ  . (9.13) This is a system of n + 6 equations in n + 6 unknowns, reflecting the fact that the floating-base system has n + 6 degrees of motion freedom. Note that a complete description of this system’s acceleration requires the two vectors a0 and ¨q. Likewise, a complete description of the velocity requires v0 and ˙q, and a complete description of the position requires the quantities 0Xref and q, where 0Xref is the Pl¨ucker transform from an inertial reference frame to floating-base coordinates.

Example 9.2 Starting with Eq. 9.13, suppose we subtract F T(Ic 0)−1 times the first row from the second. The resulting equation is Ic 0 F 0 H −F T(Ic 0)−1F  a0 ¨q  +  pc 0 C −F T(Ic 0)−1pc 0  = 0 τ  . Extracting the bottom row, we have Hfl¨q + Cfl= τ , (9.14) where Hfl= H −F T(Ic 0)−1F (9.15) and Cfl= C −F T(Ic 0)−1pc 0 . (9.16) Equation 9.14 is interesting because it presents us with a direct relationship between ¨q and τ, and the coefficients Hfland Cflcan be regarded as floating-base analogues to the coefficients of the standard equation of motion. However, this equation is not necessarily a good choice for computational purposes because the subtraction of F T(Ic 0)−1F from H in Eq. 9.15 erases any branch-induced sparsity that may have been present in H. Thus, Hflis a dense matrix; and it is perfectly possible for Hflto have more nonzero elements than H, despite being smaller. 9.4 Floating-Base Forward Dynamics If software is available to calculate fixed-base dynamics, and this software supports 6-DoF joints, then it can be used to calculate floating-base dynamics. Alternatively, if dedicated floating-base software is required, then both the articulated-body algorithm and the composite-rigid-body algorithm can be adapted to the task. Table 9.4 shows the pseudocode for a floating-base articulated-body algorithm. On comparing this algorithm with the one in Table 7.1, the following differences can be seen: v0 is now an input; pass 2 has been extended to calculate IA 0 and pA 0 ; a0 is calculated from the equation IA 0 a0 + pA 0 = 0 ; (9.17) and gravitational acceleration is added to a0 right at the end. This implementation expects the vectors f x i and ag to be supplied in floating-base coordinates, rather than reference coordinates, so that it does not need to know the value of 0Xref . The symbols 0f x i and 0ag are used accordingly. Note that 0ag is not a constant, but will vary as the floating base rotates. (0ag = 0Xref refag.) The effect of a uniform gravitational field on a free-floating rigid-body system is to modify the acceleration of every body in the system by the same

for i = 1 to NB do [XJ, Si, vJ, cJ] = jcalc(jtype(i), qi, ˙qi) iXλ(i) = XJ XT(i) if λ(i)̸ = 0 then iX0 = iXλ(i) λ(i)X0 end vi = iXλ(i) vλ(i) + vJ ci = cJ + vi × vJ IA i = Ii pA i = vi ×∗Ii vi −iX∗ 0 0f x i end IA 0 = I0 pA 0 = v0 ×∗I0 v0 −0f x 0 for i = NB to 1 do Ui = IA i Si Di = ST i Ui ... ... ui = τi −ST i pA i Ia = IA i −Ui D−1 i U T i pa = pA i + Iaci + Ui D−1 i ui IA λ(i) = IA λ(i) + λ(i)X∗ i Ia iXλ(i) pA λ(i) = pA λ(i) + λ(i)X∗ i pa end a0 = −(IA 0 )−1 pA 0 for i = 1 to NB do a′ = iXλ(i) aλ(i) + ci ¨qi = D−1 i (ui −U T i a′) ai = a′ + Si ¨qi end a0 = a0 + 0ag Table 9.4: Floating-base articulated-body dynamics algorithm amount. Thus, the relative accelerations of any two bodies are independent of gravity. The algorithm in Table 9.4 exploits this fact by calculating the motion of the system in the absence of a gravitational field, and then adding ag to a0 right at the end. If you want this algorithm to return the accelerations of other bodies in the system, then ag must also be added to these accelerations before they are returned. Table 9.5 presents floating-base versions of the recursive Newton-Euler and composite-rigid-body algorithms. Together, they calculate the coefficients of Eq. 9.13, so that it can be solved for a0 and ¨q. Note that the coefficient matrix in this equation has the same branch-induced sparsity pattern as the matrix in Eq. 9.11, so it can be factorized using the algorithms of Section 6.5. On comparing the modified Newton-Euler algorithm with the original in Table 5.1, the following differences can be seen: v0 is now an input; the term Si¨qi has been removed from the expression for ai, which is duly renamed avp i , as it accounts for only the velocity-product acceleration terms; pass 2 has been extended to calculate f0; and a line has been added to set pc 0 = f0. Once again, the vectors ag and f x i are supplied in floating-base coordinates, so that the value of 0Xref is not needed, and the symbols 0ag and 0f x i have been used in the psuedocode accordingly. Setting avp 0 = −ag at the beginning causes gravitational force terms to appear in C and pc 0, which ultimately has the effect of causing the correct gravitational acceleration term to appear in a0 when Eq. 9.13 is solved. On comparing the modified composite-rigid-body algorithm with the orig-

Calculate C and pc 0: avp 0 = −0ag for i = 1 to NB do [XJ, Si, vJ, cJ] = jcalc(jtype(i), qi, ˙qi) iXλ(i) = XJ XT(i) if λ(i)̸ = 0 then iX0 = iXλ(i) λ(i)X0 end vi = iXλ(i) vλ(i) + vJ avp i = iXλ(i) avp λ(i) + cJ + vi × vJ fi = Ii avp i + vi ×∗Ii vi −iX∗ 0 0f x i end f0 = I0 avp 0 + v0 ×∗I0 v0 −0f x 0 for i = NB to 1 do Ci = ST i fi fλ(i) = fλ(i) + λ(i)X∗ i fi end pc 0 = f0 Calculate H, F and Ic 0: H = 0 for i = 0 to NB do Ic i = Ii end for i = NB to 1 do Ic λ(i) = Ic λ(i) + λ(i)X∗ i Ic i iXλ(i) Fi = Ic i Si Hii = ST i Fi j = i while λ(j)̸ = 0 do Fi = λ(j)X∗ j Fi j = λ(j) Hij = F T i Sj Hji = HT ij end Fi = 0X∗ j Fi end Table 9.5: Floating-base versions of the recursive Newton-Euler and compositerigid-body algorithms inal in Table 6.2, the following differences can be seen: the calculation of I c i has been extended to include Ic 0; the old local variable F (which must not be confused with the new 6 × n matrix F in Eq. 9.12) has been replaced by the variables Fi, which are the block columns of F appearing in Eq. 9.12; and a line has been added after the while loop to transform each Fi to floatingbase coordinates. (Each Fi is initially expressed in body i coordinates, and is progressively transformed to floating-base coordinates.) The algorithms in these tables are not the only ways to solve the problem. Examples of some alternative approaches can be found in Hooker and Margulies (1965); Hooker (1970); Roberson and Schwertassek (1988); Wittenburg (1977). 9.5 Floating-Base Inverse Dynamics The inverse-dynamics problem for a floating-base system is really a hybriddynamics problem: the joint accelerations are known, but the base acceleration is not. Therefore, it would be possible to solve this problem using one of the general hybrid-dynamics algorithms in Sections 9.1 and 9.2. Alternatively, we could design a special-purpose algorithm, which would solve the problem more efficiently. This section presents one such algorithm. Let fi be the spatial force that is exerted on body i by its parent. In the

case of body 0, we have f0 = 0 because this body does not have a parent. In all other cases, fi is the force transmitted across joint i. In any kinematic tree, fi supports the motion of all the bodies in the set ν(i), so we have fi = X j∈ν(i) (Ij aj + vj ×∗Ij vj −f x j ) . (9.18) Our objective is to solve this equation for a0. To accomplish this, we first express ai in the form ai = a0 + ar i , where ar i is the relative acceleration of body i, and then substitute this expression into Eq. 9.18. The result is fi =  X j∈ν(i) Ij  a0 + X j∈ν(i) (Ij ar j + vj ×∗Ij vj −f x j ) . (9.19) The coefficient of a0 in this equation is the composite-rigid-body inertia Ic i . If we treat the rest of the right-hand side as a bias force, then we have fi = Ic i a0 + pc i , (9.20) where pc i = X j∈ν(i) (Ij ar j + vj ×∗Ij vj −f x j ) . (9.21) pc i is the force that would be required to support the motion of all the bodies in ν(i) if a0 happened to be zero. It can be calculated recursively using the equation pc i = pi + X j∈µ(i) pc j (9.22) where pi = Ii ar i + vi ×∗Ii vi −f x i . (9.23) When i = 0, Eq. 9.20 can be solved directly for a0, giving a0 = −(Ic 0)−1 pc 0 . (9.24) Once a0 is known, the joint forces can be calculated using τi = ST i fi = ST i (Ic i a0 + pc i) . (9.25) Table 9.6 presents the equations and pseudocode for an algorithm based on Eqs. 9.22 to 9.25 and the usual equations for body velocities, body accelerations and composite-rigid-body inertias. The equations are expressed in body coordinates, and include coordinate transforms where necessary. In line with the other floating-base algorithms, this one expects ag and f x i to be supplied

Pass 1 ar 0 = −0ag vJi = Si ˙qi vi = iXλ(i) vλ(i) + vJi ci = ˚ Si ˙qi + vi × vJi ar i = iXλ(i) aλ(i) + ci + Si ¨qi pi = Ii ar i + vi ×∗Ii vi −iX∗ 0 0f x i Pass 2 Ic i = Ii + X j∈µ(i) iX∗ j Ic j jXi pc i = pi + X j∈µ(i) iX∗ j pc j Pass 3 0a0 = −(Ic 0)−1 pc 0 ia0 = iXλ(i) λ(i)a0 τi = ST i (Ic i ia0 + pc i) ar 0 = −0ag for i = 1 to NB do [XJ, Si, vJ, cJ] = jcalc(jtype(i), qi, ˙qi) iXλ(i) = XJ XT(i) if λ(i)̸ = 0 then iX0 = iXλ(i) λ(i)X0 end vi = iXλ(i) vλ(i) + vJ ar i = iXλ(i) ar λ(i) + cJ + vi × vJ + Si ¨qi Ic i = Ii pc i = Ii ar i + vi ×∗Ii vi −iX∗ 0 0f x i end Ic 0 = I0 pc 0 = I0 ar 0 + v0 ×∗I0 v0 −0f x 0 for i = NB to 1 do Ic λ(i) = Ic λ(i) + λ(i)X∗ i Ic i iXλ(i) pc λ(i) = pc λ(i) + λ(i)X∗ i pc i end 0a0 = −(Ic 0)−1pc 0 for i = 1 to NB do ia0 = iXλ(i) λ(i)a0 τi = ST i (Ic i ia0 + pc i) end Table 9.6: Floating-base inverse dynamics equations and algorithm in floating-base coordinates, and the symbols 0ag and 0f x i have been used accordingly. This algorithm handles gravity by initializing ar 0 to −ag rather than zero. This causes gravitational force terms to appear in pc i, which in turn cause the correct gravitational acceleration to appear in a0. Alternatively, one could set ar 0 to zero at the beginning of pass 1, and then add ag to a0 at the end of pass 3. More on this topic, and on related topics, can be found in Dubowsky and Papadopoulos (1993); Fang and Pollard (2003); Longman et al. (1987); Umetani and Yoshida (1989); Vafa and Dubowsky (1990a,b); and many other papers. Example 9.3 The total spatial momentum of a rigid-body system is the sum of the momenta of its bodies. The momentum of a floating-base system is therefore given by momentum = NB X i=0 Ii vi . In the absence of external forces, this quantity is conserved. It turns out that

the momentum is also given by the expression Ic 0v0 + F ˙q. We can prove this as follows: NB X i=0 Ii vi = NB X i=0 Ii  v0 + X j∈κ(i) Sj ˙qj  = NB X i=0 Ii v0 + NB X j=1 X i∈ν(j) Ii Sj ˙qj = Ic 0 v0 + NB X j=1 Ic j Sj ˙qj = Ic 0 v0 + F ˙q . (See Eq. 4.3 for summation over ν(j).) Using a similar argument, we can show that the kinetic energy of a floating-base system is T = 1 2 NB X i=0 vT i Ii vi = 1 2  vT 0 ˙qT  Ic 0 F F T H  v0 ˙q  . Example 9.4 In a fixed-base system, the acceleration of body i is given by ai = Ji¨q + ˙Ji ˙q, where Ji is the Jacobian of body i. In a floating-base system, this becomes ai = Ji ¨q + ˙Ji ˙q + a0 . However, from Eq. 9.13 we have Ic 0 a0 + F ¨q + pc 0 = 0 . Putting these two equations together, we get ai = Ji ¨q + ˙Ji ˙q −(Ic 0)−1(F ¨q + pc 0) = J fl i ¨q + ˙Ji ˙q −(Ic 0)−1pc 0 , where J fl i = Ji −(Ic 0)−1F . The quantity J fl i is, in effect, a floating-base Jacobian, since it describes the linear dependency of ai on ¨q (cf. Umetani and Yoshida (1989)). This matrix also relates changes in ˙q to changes in vi. For example, if an impulse within the system causes ˙q to change by an amount δ ˙q, then the consequential change in vi is δvi = J fl i δ ˙q . 9.6 Gears Many mechanical systems contain gear pairs, or components that perform a similar function. From a mathematical point of view, a gear pair is an algebraic

Base Body 1 Body 2 Body 3 Body 4 1 2 3 4 ρ1 ρ2 Figure 9.2: A two-link robot arm driven by motors and gears equation involving two joint variables. Typically, one variable is a constant multiple of the other, but nonlinear relationships are also possible. Another possibility is that more than two variables can be involved. For example, a differential gear box usually has two outputs and either one or two inputs, and therefore involves either three or four variables in total. Gear constraints can be incorporated into a kinematic tree using the same technique as was used in Section 8.11 for kinematic loops; namely, to choose a vector of independent joint variables, y, and define a function, γ, that maps y to q. Let us now study this technique with the aid of an example. Figure 9.2 shows a simple two-link robot arm in which both joints are driven by electric motors via gear pairs. This system contains the following parts: a fixed base, an upper arm, a forearm, two motors, two large gear wheels and two small gear wheels. Each motor consists of two parts: a stator and a rotor. The two stators are fixed to the base and the upper arm, respectively, and the two rotors are each fixed to a small gear wheel. The two large gear wheels are fixed to the upper arm and forearm, respectively. The system therefore contains a total of four moving bodies, defined as follows: body 1: the upper arm, the first large gear wheel and the stator of the second motor; body 2: the first small gear wheel and the rotor of the first motor; body 3: the forearm and the second large gear wheel; body 4: the second small gear wheel and the rotor of the second motor. Joints 1 and 3 are therefore the shoulder and elbow joints of the robot arm, while joints 2 and 4 are the bearings inside the motors that allow the rotors to rotate

relative to the stators. The two gear pairs impose the following constraints on the joint variables: q2 = ρ1 q1 and q4 = ρ2 q3 , where ρ1 and ρ2 are gear ratios. Figure 9.2 also shows the connectivity graph of this system, and the gearing relationships between the joints. The arrows serve only to indicate the direction in which the stated gear ratio applies. Thus, the arrow labelled ρ1 points from joint 1 to joint 2 because q2 is ρ1 times q1. To incorporate the gear constraints into the equation of motion, we first define a vector of independent variables, y, and then define a function, γ, that gives q as a function of y. Let us choose q1 and q3 to be the independent variables. y is then the vector [y1 y2]T, where y1 = q1 and y2 = q3. With this choice of variables, γ is defined as follows: γ(y) =

y1 ρ1 y1 y2 ρ2 y2

. The next step is to derive the quantities G = ∂γ/∂y and g = ˙G ˙y from γ. For our example, these quantities are G = ∂γ ∂y =

1 0 ρ1 0 0 1 0 ρ2

and g = ˙G ˙y = 0 . (g would only ever be nonzero if there were a nonlinear gear pair in the system.) Finally, the equation of motion for the kinematic tree, taking into account the gear constraints, is GTH G ¨y + GT(C + Hg) = GTτ , (9.26) where H and C are the coefficients of the equation of motion in the absence of gear constraints. (See Eq. 8.45.) As the motors exert their torques about joints 2 and 4, elements τ2 and τ4 of τ will contain the motor torques, while τ1 = τ3 = 0. The inverse dynamics of a geared kinematic tree can be computed using the method described in Section 8.12. G will usually be a very sparse matrix. In fact, it will typically contain only one nonzero element per row, and many of those elements will be ones. It is therefore recommended to exploit this sparsity. By way of example, if we treat G and H as dense matrices, then the cost of calculating GTHG for a 4 × 2 matrix G is 44 multiplications and 33 additions. In contrast, if we use the actual values of G and H for the above example, then the cost is a mere

5 multiplications and 3 additions, as shown below. Obviously, the savings will be even greater with larger matrices. GTH G =  1 ρ1 0 0 0 0 1 ρ2 

H11 0 H13 H14 0 H22 0 0 H13 0 H33 0 H14 0 0 H44

1 0 ρ1 0 0 1 0 ρ2

=  H11 + ρ2 1 H22 H13 + ρ2 H14 H13 + ρ2 H14 H33 + ρ2 2 H44  . (9.27) Example 9.5 In a typical robot mechanism, the rotors of the electric motors are much smaller than the other parts of the mechanism, and therefore account for only a small fraction of the total mass. However, they also reach much higher speeds than the bodies they are driving, and can therefore exert a substantial influence on the overall dynamics of the system. The method presented in this section is the correct way to account for the dynamic effects of rotors, but there is also a quick-and-dirty method that is popular in robotics. According to this method, we dispense with the rotors and the gearboxes, and model the actuator as a fictitious electric motor that develops ρ times the torque of the real motor, runs at 1/ρ times the speed, and has ρ2 times the rotor inertia of the real motor, where ρ is the step-down gear ratio. These inertias, which are scalar constants, are simply added to the diagonal elements of H. In particular, if Irot i is the scaled rotor inertia for the motor that is driving joint i, then Irot i is added to Hii. For this to work, joint i must be a 1-DoF joint, so that Hii is a 1 × 1 matrix. This method captures the largest inertial effects of the rotors, but it ignores the smaller inertial effects and all of the gyroscopic effects. A clue as to why this method works can be seen in Eq. 9.27. If the rotors are small, then H22, H14 and H44 will be small compared with H11, H13 and H33, but ρ1 and ρ2 will be moderately large. (Gear ratios of 100 : 1 are fairly typical.) The two rotors will therefore have their biggest effect on the two diagonal elements, where they get multiplied by the squares of the gear ratios. The quick-and-dirty method also works with the articulated-body algorithm. In this case, we modify Eq. 7.44 to read Di = ST i Ui + Irot i . (9.28) Again, joint i must be a 1-DoF joint, so that Di is a 1 × 1 matrix. 9.7 Dynamic Equivalence It is perfectly possible for two different rigid-body systems to have the same equation of motion. We shall describe such systems as dynamically equivalent. An important special case occurs when the two systems differ only in their inertia parameters. An immediate consequence of this kind of equivalence is that

S I1 I2 B1 B2 f τ Figure 9.3: A two-body floating rigid-body system the inertia parameters of a rigid-body system cannot be identified from measurements of its dynamic behaviour: one must instead identify an observable subset, which are called the base inertia parameters (Khalil and Dombre, 2002). Another consequence, which is the topic of this section, is that the inertia parameters of a system model can be adjusted without altering its behaviour. If the adjustments result in fewer parameters having nonzero values, then the cost of computing the dynamics can be reduced. Consider the rigid-body system shown in Figure 9.3. It consists of two bodies, B1 and B2, having inertias of I1 and I2, respectively, which are connected by a joint having a motion subspace of S. A spatial force f acts on B1, and a generalized force τ acts at the joint. The equation of motion for this system is I1 + XT J I2 XJ XT J I2 S STI2 XJ STI2 S  a1 ¨q  + Cf Cτ  = f τ  , (9.29) where Cf Cτ  = v1 ×∗I1 v1 + XT J (I2 ˙S ˙q + v2 ×∗I2 v2) ST(I2 ˙S ˙q + v2 ×∗I2 v2)  and v2 = XJv1 + S ˙q. Except for notational differences, this equation is just an application of Eq. 9.13 to a two-body system. In formulating this equation, we have chosen the joint’s predecessor and successor frames to serve as link coordinate frames, so that 2X1 = XJ. We now modify this system by adding a spatial inertia of I∆to I1 and subtracting I∆from I2. If we want this modification not to alter the coefficients of Eq. 9.29, then I∆must meet the following conditions:3 I∆S ≡0 , I∆≡XT J I∆XJ and STv1×∗I∆v1 ≡0 . (9.30) These conditions must be met at every joint position, and for every v1 ∈M6. It therefore follows that the conditions depend only on the joint type, and that suitable values for I∆can be found without knowing anything about the system other than the joint type. 3To obtain this set of conditions, it is necessary to assume that the joint’s motion subspace, range(S), is a constant in either B1 or B2 coordinates. This assumption does not hold for all joint types.

Joint type I∆ revolute 2 6666664 I 0 0 0 −mcz 0 0 I 0 mcz 0 0 0 0 0 0 0 0 0 mcz 0 m 0 0 −mcz 0 0 0 m 0 0 0 0 0 0 m 3 7777775 prismatic 2 6666664 Ixx Ixy Ixz 0 0 0 Ixy Iyy Iyz 0 0 0 Ixz Iyz Izz 0 0 0 0 0 0 0 0 0 0 0 0 0 0 0 0 0 0 0 0 0 3 7777775 = ¯I 0 0 0  spherical 2 6666664 0 0 0 0 0 0 0 0 0 0 0 0 0 0 0 0 0 0 0 0 0 m 0 0 0 0 0 0 m 0 0 0 0 0 0 m 3 7777775 = 0 0 0 m1  Table 9.7: Values of I∆that do not alter the equation of motion Table 9.7 shows the values of I∆that meet the conditions in Eq. 9.30 for some joint types. Let us examine the physical interpretation of these results. In the case of a spherical joint, I∆is the inertia of a particle of mass m located at the joint’s rotation centre. This result says you can add particles of masses m and −m to B1 and B2, respectively, both located at the rotation centre, and because these two particles always coincide, they cancel each other exactly. In the case of a revolute joint, I∆is the inertia of a pair of particles at different positions on the joint axis (which is the z axis). m and cz are the total mass and mass centre of the pair, and I depends on the distance between them. Again, because the particles are located on the rotation axis, if we add them to B1 and add their negatives to B2 then they always coincide, and therefore always cancel. In the case of a prismatic joint, I∆is a disembodied pure rotational inertia. Such inertias are translationally invariant; so, if we add I∆to B1 and subtract it from B2, then the two will always cancel so long as B2 never rotates relative to B1. Clearly, the subtraction of a sufficiently massive particle from B2 will result in a negative mass, and the subtraction of a sufficiently large pure rotational inertia will result in a negative rotational inertia. It is therefore possible for B2 to end up with a rigid-body inertia that is not physically realizable. Nevertheless, the equation of motion for the modified system is identical to the original. If B1 and B2 are part of a larger rigid-body system, then the above results still apply—replacing a two-body subsystem with a dynamically equivalent sub-

system does not alter the equation of motion for the system as a whole. Simplification of Rigid-Body systems The results in Table 9.7 allow a dynamics simulator to simplify the inertia parameters of any given rigid-body system by adding and subtracting values of I∆that are designed to zero as many of the inertia parameters as possible. The procedure is as follows: for i = NB to 1 do I∆= best match(jtype(i), Ii) Ii = Ii −I∆ if λ(i)̸ = 0 then Iλ(i) = Iλ(i) + I∆ end end The function best match first looks in a table of I∆formulae to find an entry for the joint type jtype(i). If no entry is found, then it returns the value 06×6. However, if a formula is found, then it returns the value of I∆that maximizes the number of zero-valued inertia parameters in the expression Ii −I∆. For example, if joint i is spherical, then the returned value of I∆has the same mass as Ii, so that the mass of Ii −I∆is zero. The point of this procedure is that it produces a predictable pattern of zero-valued parameters, which can be exploited by writing code for each special case.4 Assuming that most joints are revolute, this technique can cut the cost of the recursive Newton-Euler algorithm by about 15\%, but it is less effective on forward dynamics algorithms. It is possible for a dynamics algorithm to fail on the modified system when it would have succeeded on the original. For example, if an algorithm uses the inverses of the rigid-body inertias, then it will only work if those inverses exist. If the modifications have brought the mass of body i to zero, then its modified inertia will be singular. More on this topic, and related topics, can be found in Gautier and Khalil (1990); Khalil and Kleinfinger (1987); Khalil and Dombre (2002). Augmented Bodies The above technique resembles a much older technique called augmented bodies. They both involve inertia parameters, and they both have the same objective: cost reduction. However, the augmented-body technique relies on a rewriting of the equations of motion that collects inertia parameters into constant subexpressions, which can be computed in advance. These precomputed values can be regarded as the inertia parameters of a collection of augmented bodies. 4A better strategy is to use the symbolic simplification technique described in §10.4.

Starting with a kinematic tree, augmented body i is obtained by adding one particle to rigid body i for each joint attached to that body. The locations of these particles depend on the joints, but each has a mass equal to the mass of the subtree that is connected to body i via that joint. In a variant form of augmented body, a particle is added only for each joint in µ(i). The former tends to be used with floating-base systems, and the latter with fixed-base systems. More on this topic can be found in Balafoutis et al. (1988); Balafoutis and Patel (1989, 1991); Hooker and Margulies (1965); Hooker (1970); Roberson and Schwertassek (1988); Wittenburg (1977).
